
\label{section:related work}
The literature on causal inference is abundant, and it is beyond the scope of this paper to cover it exhaustively, although \cite{Pearl_2009} is a good reference for a broad overview of this topic. Within Rubin's framework, baseline approaches consist in using treatment as a feature, or in learning two independent classifiers on the control and on the test datasets. This latter approach has the advantage of simplicity and versatility, but may lead to a selection bias. To overcome this problem, more sophisticated methods have been proposed. A first group on method consists in adaptations of classical machine learning methods. Examples are: (i) an SVM-like approach \cite{zaniewicz2013support} where two hyperplanes are introduced and properly optimized to separate class behaviors.
    (ii) a parametric Bayesian method \cite{alaa2017bayesian} for learning the treatment effects using a multi-task Gaussian process.
    (iii) several random forest-based approaches with  split criteria specifically adapted to the problem~\cite{wager2018estimation}.
Deep learning methods have also been put to good use with examples such as:
(i) a deep neural network architecture \cite{shalit2017estimating} able to learn classifiers on the test and control populations while enforcing the minimization of an integral probability metric between the distributions of these classifiers. This work builds upon \cite{johansson2016learning} where counterfactual inference has been tackled from the perspective of domain adaptation and representation learning. 
   (ii) A number of methods using neural networks have also recently emerged~\cite{yoon2018ganite,louizos2017causal}.
Interestingly, when the test and control populations have the same size, it is shown in~\cite{Jaskowski2012} that a variable change could be used, leading to the estimation of a unique probability distribution (allowing the straightforward use of classical methods).

Our approach is different from the ones above. In our case, with binary treatment and outcome, the population has a clear causal structure. We use this fact and model the population as a mixture of four causal groups. Then, from the general knowledge of the causal structure obtained by our method, we can derive the ITE and thus compare our results with ITE-specific methods.



%the authors of~\cite{shalit2017estimating} proposed a deep neural network architecture able to learn classifiers on the test and control populations while enforcing the minimization of an integral probability metric between the distributions of these classifiers. Their work builds upon \cite{johansson2016learning} where counterfactual inference has been tackled from the perspective of domain adaptation and representation learning. 
%Likewise, a SVM-based approach, developed in~\cite{zaniewicz2013support}, uses two separating hyperplanes to explicitly model the difference in class behavior between two datasets, and authors in~\cite{alaa2017bayesian} developed a parametric Bayesian method for learning the treatment effects using a multi-task Gaussian process. 
%A class transformation approach is introduced in~\cite{Jaskowski2012} to restate the problem as the estimation of a unique probability distribution. 
%Matching methods that match units in the treatment and control group with similar observable characteristics have been used with estimated propensity scores to based on weight and stratify observations~\cite{rosenbaum1987model,robins2000marginal,dehejia2002propensity,lunceford2004stratification}.
%\cite{dehejia2002propensity} focused on matching methods in detail and complement \cite{heckman1997matching} works by discussing an array of matching estimators in the context of a different dataset, while \cite{lunceford2004stratification} used methods based on weighting and stratification of observations by quantiles of estimated propensity scores.
%Several regression tree methods have been adpated to optimize for treatment effects estimation (e.g.~\cite{athey2016recursive}).
%which introduced data-driven methods that select subpopulations to estimate treatment effect heterogeneity and to test hypotheses about the differences between the effects in different subpopulations.
%Then, causal forests have been introduced in~\cite{wager2018estimation} to construct asymptotic confidence intervals.
%in complementary of \cite{foster2011subgroup} that use  regression forest to estimate the effect of covariates on outcomes in treated and control groups separately. \cite{hill2011bayesian} used the Bayesian Additive Regression Tree (BART) method of \cite{chipman2010bart} which models the inference with an unknown function approximate by a sum of binary regression trees and  product draws from the posterior using Markov chain Monte Carlo (MCMC), after specifying a prior. %It's method has the advantage to produce coherent posterior intervals estimation.
%A number of methods using neural networks have also recently emerged~\cite{yoon2018ganite,louizos2017causal}.%alaa2017deep


%\cite{athey2015machine, rzepakowski2012decision}

%Our approach tackles a more general problem than estimating ITE in the sense that we characterize the causal groups. It is related to the work in~\cite{Imbens1997} where the causal effect also derives from hidden variables. %However, our modeling is slightly different since we enforce further constraints on the assignment of the individuals to the causal populations that make use of the information we can already infer from our partial knowledge of each individual. 
%To the best of our knowledge, there is currently no work on this specific topic.




